\documentclass[aspectratio=169, 10pt]{beamer}
\usetheme{metropolis}

% --- Edinburgh colours ----------------------------------------
\definecolor{UoERed}{HTML}{7A2318}
\definecolor{UoEGold}{HTML}{B8860B}
\definecolor{CodeBG}{HTML}{F7F3ED}
\definecolor{CodeFrame}{HTML}{D5C9B1}

\setbeamercolor{palette primary}{bg=UoERed, fg=white}
\setbeamercolor{frametitle}{bg=UoERed, fg=white}
\setbeamercolor{title separator}{fg=UoEGold}
\setbeamercolor{progress bar}{fg=UoEGold, bg=UoEGold!20}
\setbeamercolor{alerted text}{fg=UoERed}
\setbeamercolor{block title}{bg=UoERed!15, fg=UoERed}
\setbeamercolor{block body}{bg=UoERed!5}

% --- Packages -------------------------------------------------
\usepackage[T1]{fontenc}
\usepackage{lmodern}
\usepackage{amsmath, amssymb, mathtools}
\usepackage{listings}
\usepackage{graphicx, booktabs}
\usepackage{tikz}
\usetikzlibrary{arrows.meta, positioning, calc, decorations.pathreplacing}
\usepackage{hyperref}
\hypersetup{colorlinks=true, linkcolor=UoERed, urlcolor=UoERed!70!black}
\setbeamertemplate{navigation symbols}{}

\lstset{
  language=Python,
  basicstyle=\ttfamily\scriptsize,
  keywordstyle=\color{UoERed}\bfseries,
  commentstyle=\color{gray}\itshape,
  stringstyle=\color{UoEGold},
  backgroundcolor=\color{CodeBG},
  frame=single,
  rulecolor=\color{CodeFrame},
  numbers=none, breaklines=true, showstringspaces=false, tabsize=4,
  xleftmargin=4pt, xrightmargin=4pt, aboveskip=6pt, belowskip=4pt,
}

\AtBeginSection[]{
  \begin{frame}
    \vfill\centering
    \usebeamerfont{title}\insertsectionhead\par
    \vfill
  \end{frame}
}

\title{Week 2 --- Deep Learning Fundamentals}
\subtitle{Gradient Descent, Backpropagation \& Training Neural Networks}
\author{Dr Juan Zurita}
\institute{Deep Learning for Macroeconomics\\
The University of Edinburgh~$\cdot$~School of Economics}
\date{}

\begin{document}
\maketitle

\begin{frame}{The Learning Problem}
Given data or a known function, find network parameters $\boldsymbol{\theta}$ that minimise a \alert{loss function}:

$$\boldsymbol{\theta}^* = \arg\min_{\boldsymbol{\theta}} \mathcal{L}(\boldsymbol{\theta})$$

\bigskip
In economics: the ``data'' is often \textbf{equilibrium conditions} rather than observed samples.
\end{frame}

\begin{frame}{Gradient Descent}
Update rule: $\boldsymbol{\theta}_{t+1} = \boldsymbol{\theta}_t - \alpha\,\nabla_{\boldsymbol{\theta}}\mathcal{L}$

\bigskip
\begin{itemize}
\item $\alpha$: learning rate --- too large diverges, too small crawls
\item \textbf{SGD}: use a random mini-batch instead of all data
\item \textbf{Adam}: adaptive learning rates per parameter --- the workhorse
\end{itemize}

\bigskip
\begin{block}{From PNM}
You already know gradient descent and Newton's method. Adam is an adaptive extension that works well for neural networks.
\end{block}
\end{frame}

\begin{frame}{Backpropagation}
\textbf{Backprop} = chain rule applied systematically through the computational graph.

\bigskip
For a composition $f = f_L \circ f_{L-1} \circ \cdots \circ f_1$:

$$\frac{\partial \mathcal{L}}{\partial \boldsymbol{\theta}_\ell} = \frac{\partial \mathcal{L}}{\partial f_L} \cdot \frac{\partial f_L}{\partial f_{L-1}} \cdots \frac{\partial f_{\ell+1}}{\partial f_\ell} \cdot \frac{\partial f_\ell}{\partial \boldsymbol{\theta}_\ell}$$

\bigskip
PyTorch computes this automatically via \texttt{loss.backward()}.
\end{frame}

\begin{frame}{Activation Functions}
\begin{center}
\begin{tabular}{lll}
\toprule
Name & Formula & Use case \\
\midrule
ReLU & $\max(0, x)$ & Default hidden layers \\
Sigmoid & $1/(1+e^{-x})$ & Output in $[0,1]$ \\
Tanh & $(e^x - e^{-x})/(e^x + e^{-x})$ & Output in $[-1,1]$ \\
Softplus & $\ln(1 + e^x)$ & Smooth ReLU \\
\bottomrule
\end{tabular}
\end{center}

\bigskip
For economics: \textbf{softplus} is useful when outputs must be positive (e.g.\ consumption).
\end{frame}

\begin{frame}{Overfitting and Regularisation}
\begin{itemize}
\item \textbf{Overfitting}: the network memorises training points but generalises poorly
\item \textbf{Weight decay} ($L_2$ regularisation): penalise large weights
\item \textbf{Dropout}: randomly zero out neurons during training
\item \textbf{Early stopping}: halt when validation loss stops improving
\end{itemize}

\bigskip
In economics, overfitting is less of a concern when the ``data'' comes from equilibrium conditions --- but training stability still matters.

\bigskip
\textbf{Next week:} automatic differentiation --- the engine behind backpropagation.
\end{frame}

\end{document}
